\startsection[title={Fermentador}]
	O fermentador, onde ocorrem as reações bioquímicas de interesse, consiste em um balão volumétrico de três vias apoiado sobre um aquecedor e submerso em um banho de água. Um dos bocais é utilizado para a entrada de nitrogênio, enquanto outro permite a saída dos produtos gasosos. Essa saída está conectada a uma proveta invertida, preenchida com água e também submersa em água, formando uma câmara de gás à pressão aproximadamente constante.

	\startplacefigure[title={Ilustração do fermentador}, reference={fig:fermentator}]
		\externalfigure[image/fermentator.png][height=400pt]
	\stopplacefigure

	O lodo de gelatina constitui a biomassa responsável pela produção preferencial de hidrogênio e ácidos graxos voláteis a partir da fermentação de açúcar demerara. A \in{Figura}[fig:fermentator] apresenta os componentes utilizados para a contenção da cultura, do meio reacional e dos produtos, sendo a alimentação realizada apenas no início do processo. A atmosfera do reator é composta exclusivamente por nitrogênio, caracterizando um ambiente anaeróbio e quimicamente inerte. Dessa forma, os compostos disponíveis para metabolização pela cultura são intencionalmente limitados, tornando o substrato a principal fonte de carbono e energia para a produção dos metabólitos de interesse.
\stopsection

